\begin{abstract}

Large Vision-Language Models (LVLMs) suffer from \textit{hallucination} issues, wherein the models generate plausible-sounding but factually incorrect outputs, undermining their \textit{reliability}. A comprehensive quantitative evaluation is necessary to identify and understand the extent of hallucinations in these models. However, existing benchmarks are often limited in scope, focusing mainly on object hallucinations. Furthermore, current evaluation methods struggle to effectively address the subtle semantic distinctions between model outputs and reference data, as well as the balance between hallucination and informativeness. To address these issues, we introduce a multi-dimensional benchmark covering objects, attributes, and relations, with challenging images selected based on associative biases. Moreover, we propose a large language model (LLM)-based two-stage evaluation framework that generalizes the popular CHAIR metric~\cite{rohrbach-etal-2018-object} and incorporates both faithfulness and coverage into the evaluation. Experiments on 10 established LVLMs demonstrate that our evaluation metric is more comprehensive and better correlated with humans than existing work when evaluating on our challenging human-annotated benchmark dataset. Our work also highlights the critical balance between \textit{faithfulness} and \textit{coverage} of model outputs, and encourages future works to address hallucinations in LVLMs while keeping their outputs \textit{informative}.\footnote{Our dataset and code can be found here: \url{https://github.com/haoyiq114/VALOR}.}
   
\end{abstract}